\scanpage{32}
% L1–5, upper paragraph
1°) Le groupe des transformations infinitésimales \inserted{de type $S'/S$} de $X_S/Y_S$ \uline{opère} sur $\Gamma(X_{S'}/Y_{S'})$ des sections \unsure{proches} de type $S'$ des sections de $X/Y$, en induisant l'identité sur le quotient $\Gamma(X/Y)$ de ce dernier ensemble, de plus

% L6–10, formula with gloss running below the brace
\[
\Gamma(X/Y)\longrightarrow
\Gamma\underbrace{(X_{S'}/Y_{S'})}_{\substack{\text{ensemble}\\
\text{: groupe}\\ \text{d'opérateurs}\\
\operatorname{Aut}^{\{\Lambda\}}_{S}(X_{S'},Y_{S'})}}
\simeq\Gamma(X_Y^{\{Y_{S'}\}}/Y)
\]

% L11–17, central note
\uline{N.B.} $S,S'$ n'interviennent pas effectivement \deleted{\unsure{mais}} \inserted{[\unsure{seult|seulement} par} $X_{S'}\to Y_{S'}$] dans la définition \deleted{\unsure{de}} la première ligne, mais seulement dans la définition du groupe d'opérateurs sur $\Gamma(X_{S'}/Y_{S'})$, : les transformations inf. de type $S'$ de $X_{\deleted{\unsure{S}}S}/Y_S$ i.e. les $S'$-automorphismes de $Y_{S'}$ et $X_{S'}$ .....

% L18–20, lower paragraph
2°) Fixons un $w$. Un $Y$-hom $\displaystyle X\xrightarrow{h}X'$ est compatible avec des connexions $w=(u,v)$, $w'=(u',v)$ si on a commutativité \unsure{du|de}

% L21–22, diagram, and incomplete continuation at right
\diagram{\begin{tikzcd}[ampersand replacement=\&]
X_{S'}\arrow[r,"h_{S'}"]\arrow[d,"u"'] \& X'_{S'}\arrow[d,"u'"]\\
X_{S'}\arrow[r,"h_{S'}"] \& X'_{S'}
\end{tikzcd}}
\hfill i.e. $\deleted{\unsure{X'|h'}}$ \unsure{l}

% archive sheet number: 31

